% Section 1: Introduction
\subsection{Introduction}
\label{sec:forecasting_rl_intro}

Forecasting and reinforcement learning have evolved on adjacent benches. Time series
forecasters minimise mean squared error or a strictly proper scoring rule.
Reinforcement learners maximise expected discounted reward. The two fields share an
underlying object, namely the conditional distribution of a future state, yet they
disagree on which functional of that distribution counts as evaluation.
\citet{GrangerPesaran2000} stated the problem two decades ago in a sentence that has
not been fully absorbed in econometrics. ``The literature on forecast evaluation in
econometrics has been primarily concerned with statistical accuracy measures of point
forecasts, and little attention has been paid to the economic importance of these
forecasts.'' Forecasts are not made in isolation, they are designed for use in a
decision, and the decision is rarely a point estimate of conditional expectation.

The classical argument is easiest to state in the two-state two-action setting of
\citet{GrangerPesaran2000}. Let $y_t \in \{0, 1\}$ index a Bad and a Good state and
let $a_t \in \{0, 1\}$ index a No and a Yes action. The four-cell payoff matrix
$\{U_{yg}, U_{nb}, U_{yb}, U_{ng}\}$ defines a cost-benefit ratio
$c_t = (U_{nb} - U_{yb}) / (U_{yg} - U_{ng} + U_{nb} - U_{yb})$ and the optimal
decision is to act if and only if the forecast probability $\hat\pi_t$ of the Bad
event exceeds the threshold $\theta_t = c_t / (1 + c_t)$, namely
$\hat a_t = \mathbf 1\{\hat\pi_t > \theta_t\}$. The economic value of using the
forecast $\hat\pi_t$ to decide, taken in expectation over the data generating process,
is maximised at $\hat\pi_t = \pi_t$, the true conditional probability. This is the
\emph{supreme solution} of \citet{GrangerPesaran2000} and it has the convenient
property that it does not depend on the payoff ratio $\theta_t$, so the same forecast
serves all decision makers regardless of their cost structure.

Two scoring devices arise as soon as one tries to evaluate a forecast against this
benchmark. The Brier score \citep{Brier1950}
$B = T^{-1} \sum_{t=1}^T (y_t - \hat\pi_t)^2$ is the cross-sectional analogue of
mean squared error and admits the decomposition of \citet{Murphy1973} into
reliability, resolution, and uncertainty. The Kuipers score
$\mathrm{KS} = H(\theta) - F(\theta)$, where $H$ and $F$ are the hit rate and the
false alarm rate at threshold $\theta$, is the directional analogue used in
meteorology and is shown by \citet{GrangerPesaran2000} to coincide with a
standardised version of the \citet{PesaranTimmermann1992} market-timing test up to
the variance scaling, $\hat\pi - \tilde\pi^\star = 2 P_B (1 - P_B) \,\mathrm{KS}$.
The relationship between the average economic value and the Kuipers score is exact
only in the knife-edge case $\theta_t = \bar y$, namely when the payoff ratio
coincides with the unconditional event frequency. Outside this case a forecast that
ranks higher on Brier or on Kuipers can rank lower in realised economic value, and
two forecasts with identical statistical accuracy can be ordered in opposite
directions by the relevant decision loss.

\citet{GneitingRaftery2007} formalised the statistical complement of the argument.
A scoring rule $S(P, y)$ is proper, with respect to a class of distributions
$\mathcal P$, if for every $P, Q \in \mathcal P$ the expected score satisfies
$\mathbb E_{y \sim Q}[S(Q, y)] \ge \mathbb E_{y \sim Q}[S(P, y)]$, and it is strictly
proper if the inequality is strict for $P \neq Q$. Their Theorem 1 establishes that
every regular proper scoring rule arises from a convex generating function $G$ on
$\mathcal P$ via $S(P, \omega) = G(P) + G^\star(P, \omega) - \int G^\star(P, \cdot)\,
\mathrm dP$, where $G^\star$ is a subtangent of $G$ at $P$. The Brier score
corresponds to the convex function $G(p) = \|p\|_2^2$ on the simplex, the logarithmic
score to negative Shannon entropy, and the continuous ranked probability score and
the energy score follow from generalisations of the same construction to richer
sample spaces. The decision-theoretic interpretation due to \citet{Dawid1998} closes
the loop. Given an outcome space $\Omega$, an action space $\mathcal A$, a utility
$U(\omega, a)$, and the Bayes act $a_P = \arg\max_a \mathbb E_{P}[U(\omega, a)]$, the
score $S(P, \omega) = U(\omega, a_P)$ is proper with respect to $\mathcal P$, with
strict propriety holding when the optimal Bayes act is unique for each forecast.

The two halves of the present chapter sit on either side of this observation.
Reinforcement learning has, since the early 2000s, been brought \emph{to} forecasting
in four interlocking ways. RL is used as a meta-algorithm that selects ensemble
weights for a portfolio of base forecasters, as a mechanism for training a
probabilistic forecaster directly on a proper scoring rule that doubles as the
downstream decision loss, as the unifying frame for approximate dynamic programming
in sequential estimation, and most consequentially as the engine behind
decision-focused learning, where the predictor is trained on an estimate of the
realised decision regret rather than on a prediction-accuracy loss. The
methodological frontier here is \citet{ElmachtoubGrigas2022} and
\citet{DontiAmosKolter2017}. Conversely, forecasting has been brought \emph{to} RL
through world models, trajectory generative models, counterfactual outcome predictors
for off-policy evaluation, and the recent generation of pretrained time series
foundation models. The methodological frontier here is the Dreamer family of
\citet{Hafner2025} on the world-model side and the Chronos line of \citet{Ansari2024}
on the foundation-model side.

The bridge between the two halves is the question of which loss should govern the
training of the forecaster when the forecast is destined for use inside a decision.
\citet{ElmachtoubGrigas2022} and \citet{DontiAmosKolter2017} can be read as Wald
(1950) statistical decision theory made differentiable. Wald defined the loss of a
decision rule as the expected payoff under the true distribution and the optimal
procedure as the minimax rule against the worst admissible distribution. SPO$^+$
and task-based learning operationalise this principle by parametrising the
predictor, defining a differentiable estimate of the decision-induced loss, and
fitting the parameters by stochastic gradient descent. The risk bounds of
\citet{LiuGrigas2021} and \citet{HuangGupta2024} turn the qualitative argument into
a quantitative statement about sample complexity. The picture across these strands
is that the Granger-Pesaran question has a constructive answer at scale, conditional
on misspecification and conditional on the downstream decision being expressible in
a differentiable way.

We are economists, so we frame each piece by what it adds beyond ordinary least
squares with a plug-in decision rule. The chapter is organised in seven sections.
Section~\ref{sec:rl_for_forecasting} covers RL for forecasting, starting with the
shallow uses (RL as an ensemble combiner, RL as a vehicle for Brier-loss reward
training) and ending with the deepest use (decision-focused learning).
Section~\ref{sec:forecasting_for_rl} covers forecasting for RL, organised around
world models, trajectory generative models, counterfactual forecasting for
off-policy evaluation, and time series foundation models.
Section~\ref{sec:bridge} returns to the loss question, with explicit attention to
the three obstacles that prevent the immediate replacement of standalone forecasting
by decision-focused training. Section~\ref{sec:forecasting_rl_applications} surveys
four application domains where the prediction-and-decision interaction is the active
research question. Section~\ref{sec:forecasting_rl_sim} is a synthetic simulation on
AR(1) demand under heteroscedastic misspecification, designed to make the central
methodological point concrete on a familiar substrate. Section~\ref{sec:forecasting_rl_open}
closes with four problems an economist is well placed to attack.
